\documentclass[a4paper, 12pt]{article}
\usepackage[hagavi]{azzam}

\title{Exam 2}
\date{Sabtu, 25 Juli 2026}
\begin{document}

\maketitle

\begin{enumerate}
    \item Dalam sebuah kelompok yang terdiri dari $20200$ pria dan $2020$ wanita, akan dipilih sepasang orang secara acak. Pemilihan dilakukan terus-menerus serta orang yang sudah dipilih dikeluarkan dari kelompok dan tidak akan dipilih lagi di pemilihan selanjutnya. Tentukan banyaknya minimal pemilihan yang dilakukan sedemikian sehingga pasti terdapat sepasang orang yang terdiri dari dua pria.
    \item Misalkan $f:\mathbb{N}\rightarrow\mathbb{Z}$ adalah sebuah fungsi dengan $f(1)=0$, $f(2)=1$, $f(3n+1)=f(3n)=f(n)$ dan $f(3n+2)=f(n)+1$. Tentukan banyaknya bilangan asli $k<2020$ sedemikian sehingga $f(k)=f(2020)$.
    \item Tentukan, dengan bukti, nilai dari
    \[ \sum_{b=5}^{\infty}\frac{15b^2+68b+120}{2^b(b^4+64)}. \]

    \item Diberikan segitiga lancip $ABC$ dengan $AB \neq AC$. Definisikan titik $I$ sebagai titik pusat lingkaran dalam segitiga $ABC$ serta titik $I_A$ sebagai titik pusat lingkaran singgung luar segitiga $ABC$ yang berhadapan dengan titik $A$. Selain itu, definisikan pula titik $X$ dan $Y$ sebagai titik pada lingkaran luar segitiga $ABC$ sedemikian sehingga 
    \[ \angle AXI = \angle AYI_A = 90^\circ. \]
    Apabila $XI$ menyinggung lingkaran luar segitiga $AYI$, buktikan bahwa $XY$ menyinggung lingkaran dalam segitiga $ABC$.
\end{enumerate}


\end{document}